\documentclass[
  12pt,
  a4paper,
  twoside,
  brazilian
]{article}

\usepackage[utf8]{inputenc}
\usepackage[T1]{fontenc}
\usepackage{times}
\usepackage[brazilian]{babel}

\usepackage[
  top=3cm, left=3cm, right=2cm, bottom=2cm
]{geometry}

\usepackage{setspace}
\singlespacing

\usepackage{indentfirst}
\setlength{\parindent}{1.25cm}

\usepackage{fancyhdr}
\setlength{\headheight}{14.5pt}
\pagestyle{fancy}
\fancyhf{}
\fancyhead[LE,RO]{\thepage}
\fancyhead[RE,LO]{\textit{Raciocínios Científicos em Exames de Pós-Graduação}}
\renewcommand{\headrulewidth}{0.5pt}

\usepackage[colorlinks=true, linkcolor=black, citecolor=black, urlcolor=black]{hyperref}
\usepackage{array, booktabs, longtable}
\usepackage{graphicx}
\usepackage[bottom]{footmisc}

\usepackage{titlesec}
\titleformat{\section}{\normalfont\bfseries\large}{\thesection}{0.5em}{}
\titleformat{\subsection}{\normalfont\bfseries\normalsize}{\thesubsection}{0.5em}{}
\titleformat{\subsubsection}{\normalfont\bfseries\normalsize}{\thesubsubsection}{0.5em}{}

\usepackage{tocloft}
\renewcommand{\cftsecleader}{\cftdotfill{\cftdotsep}}
\providecommand{\tightlist}{%
  \setlength{\itemsep}{0pt}\setlength{\parskip}{0pt}}
\newcommand{\tabnota}[1]{\vspace{2mm}\par\noindent{\footnotesize #1}}

\begin{document}

% --- Folha de rosto (com identificação do autor) ---
% Esta página é removida para a avaliação cega (blind review)
\pagenumbering{Alph}
\begin{titlepage}
  \begin{center}
    \vspace*{3cm}
    {\large\bfseries Marcelo Claro Laranjeira\par}
    \vspace{0.5cm}
    {\normalsize \lbrack Mestrando em Educação para a Ciência\rbrack, \lbrack Função\rbrack,
    \lbrack Departamento\rbrack, \lbrack Unidade\rbrack, \\
    \lbrack Universidade por extenso (sigla)\rbrack. \lbrack Cidade\rbrack, \lbrack UF\rbrack. \\
    \textit{\lbrack e-mail@exemplo.com\rbrack}\par}
    \vspace{3cm}
    {\LARGE\bfseries Raciocínios Científicos em Exames de Seleção para Pós-Graduação: Uma Análise Comparativa entre Física e Matemática à Luz da Taxonomia 350 e da Taxonomia de Bloom\par}
    \vspace{0.5cm}
    {\large\bfseries Scientific Reasoning in Graduate Admission Exams: A Comparative Analysis between Physics and Mathematics through Taxonomy 350 and Bloom's Taxonomy\par}
    \vfill
    {\large 2026\par}
  \end{center}
\end{titlepage}

\clearpage
\pagenumbering{arabic}
\tableofcontents
\newpage

\section{Raciocínios Científicos em Exames de Seleção para
Pós-Graduação: Uma Análise Comparativa entre Física e Matemática à Luz
da Taxonomia 350 e da Taxonomia de
Bloom}\label{raciocuxednios-cientuxedficos-em-exames-de-seleuxe7uxe3o-para-puxf3s-graduauxe7uxe3o-uma-anuxe1lise-comparativa-entre-fuxedsica-e-matemuxe1tica-uxe0-luz-da-taxonomia-350-e-da-taxonomia-de-bloom}

\subsection{Scientific Reasoning in Graduate Admission Exams: A
Comparative Analysis between Physics and Mathematics through Taxonomy
350 and Bloom's
Taxonomy}\label{scientific-reasoning-in-graduate-admission-exams-a-comparative-analysis-between-physics-and-mathematics-through-taxonomy-350-and-blooms-taxonomy}

\textbf{Resumo} Este artigo analisa comparativamente 150 questões de exames de seleção para pós-graduação -- 60 do Exame Unificado de Física (EUF) e 90 do Exame Nacional de Acesso de Matemática (ENA) -- classificando-as segundo a Taxonomia 350 de Raciocínios Científicos e a Taxonomia de Bloom revisada. Os resultados revelam padrões taxonomicamente distintos: 86,7\% das questões de Física concentram-se na categoria XIII (Físico-Matemáticos), enquanto 67,8\% das questões de Matemática pertencem à categoria I (Lógico-Dedutivos). O teste exato de Fisher indica odds ratio de 61,0 (p < 0,001) para a associação entre Matemática e raciocínio lógico-dedutivo. Na dimensão de Bloom, 58\% das questões exigem Aplicar e 42\% Analisar, com distribuição homogênea entre os exames (V de Cramer = 0,047; p = 0,566). Foram identificados 25 R-codes distintos, com R008 (Dedução Formal Passo-a-Passo) concentrando 50\% das questões de Matemática. A análise evidencia que os exames avaliam perfis epistemológicos distintos: a Física privilegia raciocínios físico-matemáticos e de modelagem, enquanto a Matemática enfatiza o raciocínio lógico-dedutivo formal. Discutem-se implicações para o desenho de instrumentos de seleção e para o ensino de ciências e matemática na pós-graduação.

\textbf{Palavras-chave:} Taxonomia 350. Raciocínio científico. Taxonomia
de Bloom. EUF. ENA. Seleção pós-graduação.

\textbf{Abstract} \emph{This paper comparatively analyzes 150 questions
from graduate admission exams --- 60 from the Unified Physics Exam (EUF)
and 90 from the National Mathematics Access Exam (ENA) --- classifying
them according to the Taxonomy 350 of Scientific Reasoning and the
Revised Bloom's Taxonomy. Results reveal taxonomically distinct
patterns: 86.7\% of Physics questions concentrate in category XIII
(Physico-Mathematical), while 67.8\% of Mathematics questions belong to
category I (Logical-Deductive). Fisher's exact test yields an odds ratio
of 61.0 (p \textless{} 0.001) for the association between Mathematics
and logical-deductive reasoning. In the Bloom dimension, 58\% of
questions require Applying and 42\% Analyzing, with homogeneous
distribution across exams (Cramer's V = 0.047; p = 0.566). Twenty-five
distinct R-codes were identified, with R008 (Step-by-Step Formal
Deduction) concentrating 50\% of Mathematics questions. The analysis
shows that the exams assess distinct epistemological profiles: Physics
privileges physico-mathematical reasoning and modeling, while
Mathematics emphasizes formal logical-deductive reasoning. Implications
for admission instrument design and for science and mathematics
education in graduate programs are discussed.}

\textbf{Keywords:} Taxonomy 350. Scientific reasoning. Bloom's taxonomy.
EUF. ENA. Graduate admission.

\begin{center}\rule{0.5\linewidth}{0.5pt}\end{center}

\section{Introdução}\label{introduuxe7uxe3o}

A seleção de candidatos para programas de pós-graduação em ciências
exatas no Brasil utiliza, com frequência, exames escritos como principal
instrumento de avaliação. Essa prática, consolidada ao longo de décadas,
baseia-se na premissa de que o desempenho em uma prova de conhecimentos
específicos constitui indicador razoável da capacidade do candidato para
realizar pesquisa em nível de mestrado ou doutorado. No caso da Física,
o Exame Unificado de Física (EUF) é aplicado anualmente por mais de
trinta programas associados, abrangendo conteúdos de seis áreas:
Mecânica Clássica, Eletromagnetismo, Termodinâmica, Mecânica
Estatística, Mecânica Quântica e Física Moderna. Cada edição do EUF
contém 10 questões, distribuídas de forma a equilibrar a representação
das seis áreas. Para a Matemática, o Exame Nacional de Acesso (ENA)
seleciona candidatos para o Mestrado Profissional em Matemática em Rede
Nacional (Profmat), programa que atende a centenas de professores de
Matemática em todo o território nacional. O ENA contém 30 questões por
edição, cobrindo áreas como Álgebra, Geometria, Teoria dos Números,
Análise e Estatística.

A despeito da relevância desses exames para a formação de pesquisadores
e professores no país, pouco se conhece sobre a natureza dos processos
cognitivos e dos raciocínios científicos que eles efetivamente avaliam.
Em outras palavras, sabe-se quais conteúdos são cobrados, mas não se
sabe com clareza que tipo de raciocínio é exigido para resolvê-los, nem
se esse perfil é coerente com as competências almejadas pelos programas
de pós-graduação. Essa lacuna é relevante porque diferentes tipos de
raciocínio científico podem estar associados a diferentes desempenhos
acadêmicos e a diferentes perfis de sucesso na pesquisa. Um exame que
privilegia exclusivamente a dedução formal pode, por exemplo, selecionar
candidatos hábeis em manipulação simbólica, mas sub-representar aqueles
com talento para modelagem de fenômenos ou para o raciocínio
probabilístico. Estudos prévios em educação em ciências têm se dedicado
a analisar exames de larga escala como o ENEM (Silva; Dias, 2023)\footnote{\url{https://doi.org/10.55767/2451.6007.v33.n2.35304}}, o
ENADE (Costa, 2017)\footnote{\url{https://doi.org/10.5007/2175-7941.2017v34n3p697}} e exames internacionais como o TIMSS (Mullis;
Martin, 2017), empregando predominantemente a Taxonomia de Bloom (Bloom
\textit{et al.}, 1956; Krathwohl, 2002)\footnote{\url{https://doi.org/10.1207/s15430421tip4104_2}} como quadro de referência. Essa abordagem
tem se mostrado produtiva para descrever o perfil geral de complexidade
cognitiva das questões. Contudo, a Taxonomia de Bloom, embora útil para
classificar níveis de complexidade cognitiva, não captura a
especificidade dos raciocínios científicos mobilizados em cada questão
--- isto é, não distingue, por exemplo, entre um raciocínio dedutivo, um
raciocínio por modelagem matemática ou um raciocínio probabilístico,
todos os quais podem situar-se no mesmo nível de Aplicar ou Analisar.

Para suprir essa lacuna, Padilha (2020) propôs a Taxonomia 350 de
Raciocínios Científicos, um sistema de classificação que organiza 350
tipos de raciocínio científico (R-codes) em categorias epistemológicas
mais amplas. Essa taxonomia permite identificar, com granularidade fina,
quais habilidades de raciocínio são exigidas por cada questão, indo além
dos níveis genéricos de Bloom. A Taxonomia 350 organiza-se em categorias
que vão desde raciocínios lógico-dedutivos (Categoria I) até raciocínios
quânticos e relativísticos (Categoria XV), passando por raciocínios
analítico-estruturais, probabilísticos, físico-matemáticos, entre
outros.

O presente estudo situa-se na interface entre essas duas tradições
taxonômicas, buscando integrá-las em uma análise articulada e
multidimensional. Seu objetivo é analisar e comparar 150 questões de
exames de seleção para pós-graduação --- 60 do EUF (Física) e 90 do ENA
(Matemática) --- classificando-as simultaneamente segundo a Taxonomia
350 e a Taxonomia de Bloom revisada (Krathwohl, 2002). As perguntas que
orientam a investigação são três: (i) Quais raciocínios científicos
(R-codes e categorias) são predominantes em cada exame? (ii) Como se
distribuem os níveis cognitivos de Bloom entre as questões de Física e
Matemática? (iii) Existem associações estatisticamente significativas
entre o tipo de exame, a categoria taxonômica e o nível de Bloom?

Estudos comparativos entre disciplinas no âmbito da classificação
taxonômica de exames são escassos na literatura brasileira. Trabalhos
como o de Lithner (2008)\footnote{\url{https://doi.org/10.1007/s10649-007-9066-2}} sobre raciocínio matemático em exames
universitários suecos e o de Teodorescu \textit{et al.}~(2013)\footnote{\url{https://doi.org/10.1103/PhysRevSTPER.9.010105}} com a taxonomia
TIPP para problemas de Física oferecem quadros de referência relevantes,
mas não realizam a comparação direta entre as duas disciplinas no
contexto da pós-graduação brasileira. A presente pesquisa busca
preencher esse vazio, oferecendo subsídios tanto para a construção de
instrumentos de seleção quanto para a reflexão sobre o perfil
epistemológico que cada exame privilegia.

O artigo está organizado em seis seções. A seção 2 apresenta o
referencial teórico, abrangendo a Taxonomia de Bloom revisada e a
Taxonomia 350. A seção 3 descreve a metodologia, incluindo o corpus de
análise, o procedimento de classificação e os métodos estatísticos
empregados. A seção 4 expõe os resultados organizados por dimensão de
análise. A seção 5 discute os resultados à luz da literatura, apontando
implicações e limitações. Por fim, a seção 6 tece as considerações
finais e sugere direções para pesquisas futuras.

\newpage
\section{Referencial Teórico}\label{referencial-teuxf3rico}

\subsection{A Taxonomia de Bloom e sua
Revisão}\label{a-taxonomia-de-bloom-e-sua-revisuxe3o}

A Taxonomia de Bloom, originalmente publicada em 1956 por Benjamin Bloom
e colaboradores, estabeleceu uma hierarquia de objetivos educacionais
organizada em seis níveis cognitivos: Conhecimento, Compreensão,
Aplicação, Análise, Síntese e Avaliação (Bloom \textit{et al.}, 1956). Essa
taxonomia tornou-se um dos referenciais mais influentes na educação,
sendo amplamente utilizada para planejamento curricular, elaboração de
objetivos de aprendizagem e análise de instrumentos avaliativos (Ferraz;
Belhot, 2010). Sua popularidade decorre, em grande medida, da
simplicidade intuitiva da hierarquia proposta: níveis mais elementares
(Conhecimento, Compreensão) envolvem processos cognitivos básicos,
enquanto níveis mais elevados (Síntese, Avaliação) requerem operações
mentais mais complexas e integradas.

Em 2001, Anderson e Krathwohl lideraram uma revisão substantiva da
taxonomia original, resultando em duas dimensões independentes: a
dimensão do Conhecimento (Factual, Conceitual, Procedimental e
Metacognitivo) e a dimensão dos Processos Cognitivos (Lembrar, Entender,
Aplicar, Analisar, Avaliar e Criar) (Krathwohl, 2002). Essa revisão
substituiu os substantivos da versão original por verbos --- por
exemplo, ``Aplicação'' tornou-se ``Aplicar'' --- e reorganizou a
hierarquia, posicionando ``Criar'' como o nível mais complexo, acima de
``Avaliar''.

No contexto brasileiro, Ferraz e Belhot (2010)\footnote{\url{https://doi.org/10.1590/S0104-530X2010000200015}} realizaram um dos
trabalhos mais citados sobre a aplicação da taxonomia revisada no ensino
superior, discutindo sua utilidade para alinhar objetivos, atividades e
avaliações. Os autores argumentam que a principal contribuição da
taxonomia revisada é fornecer uma linguagem comum para educadores de
diferentes áreas, permitindo que objetivos instrucionais sejam
comunicados de forma precisa e que avaliações sejam desenhadas para
capturar níveis específicos de desempenho cognitivo. Na área de ensino
de Física, Silva e Dias (2023) aplicaram a taxonomia para analisar 86
questões de Física do ENEM entre 2014 e 2019, constatando que 52,3\% das
questões concentravam-se nos níveis de Aplicar e Analisar. Esse
resultado é coerente com a natureza do ENEM, que busca avaliar
competências aplicadas em contextos do mundo real, mas também levanta
questões sobre a representatividade dos níveis mais elevados (Avaliar,
Criar) em exames de larga escala. Em Matemática, estudos como o de Costa
(2017) sobre o ENADE também evidenciaram a predominância dos níveis
intermediários da taxonomia, com concentração similar nos níveis de
Aplicar e Analisar.

A presente pesquisa adota a versão revisada de Krathwohl (2002) e
concentra-se nos dois níveis mais frequentes em exames de seleção ---
Aplicar e Analisar ---, uma vez que exames que demandam exclusivamente
Lembrar ou Entender seriam insuficientes para selecionar candidatos à
pós-graduação, e itens nos níveis Avaliar e Criar são raros em provas
escritas de múltipla escolha ou dissertativas de tempo controlado.

\subsection{A Taxonomia 350 de Raciocínios
Científicos}\label{a-taxonomia-350-de-raciocuxednios-cientuxedficos}

A Taxonomia 350, desenvolvida por Padilha (2020) no âmbito do projeto
``Estilos de Raciocínio Científico'' do Instituto de Estudos Avançados
da USP, constitui um sistema de classificação de raciocínios científicos
com 350 tipos distribuídos em quinze categorias epistemológicas.
Diferentemente da Taxonomia de Bloom, que classifica a
\emph{complexidade} do processo cognitivo, a Taxonomia 350 classifica a
\emph{natureza} do raciocínio científico mobilizado.

As quinze categorias da Taxonomia 350 abrangem desde raciocínios
Lógico-Dedutivos (Categoria I) --- que incluem a dedução formal
passo-a-passo (R008), a indução estrutural (R011), a redução ao absurdo
(R023) e a crítica por contraexemplo (R009) --- até raciocínios mais
especializados como os Físico-Matemáticos (Categoria XIII), que agrupam
a modelagem matemática (R183), a dualidade e transformação (R071), o
princípio variacional (R069), a simetria e conservação (R065) e a
análise por Monte Carlo (R117), entre outros. Entre essas categorias,
encontram-se também raciocínios Analítico-Estruturais (Categoria II),
que envolvem a decomposição de problemas complexos em partes
gerenciáveis; raciocínios Probabilísticos e Estatísticos (Categoria IV),
fundamentais para o tratamento de incertezas; raciocínios
Econômico-Decisórios (Categoria VII), que envolvem otimização e análise
de custo-benefício; e raciocínios de Lógica Matemática Avançada
(Categoria XIV), que incluem raciocínios sobre propriedades abstratas de
estruturas matemáticas. A Categoria II (Analítico-Estruturais) contempla
a decomposição modular (R010), enquanto a Categoria VII
(Econômico-Decisórios) abriga a otimização restrita (R054) e o
raciocínio marginal-incremental (R050). A Categoria XIV (Lógica
Matemática Avançada) inclui o raciocínio espectral de operadores (R109)
e a homotopia fundamental (R229).

Cada R-code é definido operacionalmente por uma descrição do tipo de
raciocínio, acompanhada de exemplos de aplicação em contextos
científicos. A taxonomia permite, assim, que um analista classifique uma
questão não apenas pelo ``quão complexa'' ela é, mas pelo ``que tipo de
raciocínio'' ela demanda --- uma distinção importante para a compreensão
dos perfis cognitivos avaliados em exames de seleção.

\subsection{Estudos de Classificação Taxonômica em
Exames}\label{estudos-de-classificauxe7uxe3o-taxonuxf4mica-em-exames}

A literatura internacional registra diversos esforços de classificação
taxonômica de exames. Lithner (2008) propôs um arcabouço para analisar o
raciocínio matemático em exames universitários, distinguindo entre
raciocínio imitativo (baseado em memorização ou algoritmos) e raciocínio
criativo (fundamentado em argumentação matemática). Em uma análise de
593 tarefas de exames suecos, o autor constatou que mais de 80\%
demandavam apenas raciocínio imitativo, levantando questões sobre o
alinhamento entre os objetivos curriculares e a avaliação efetiva.

Teodorescu \textit{et al.}~(2013) desenvolveram a TIPP (Taxonomy of Introductory
Physics Problems), uma taxonomia específica para problemas introdutórios
de Física que considera três dimensões: tipo de problema, quantidade de
física envolvida e nível cognitivo. Em uma aplicação com 450 problemas,
os autores demonstraram que a taxonomia permitia distinguir problemas
rotineiros de problemas que exigiam integração conceitual mais profunda.

No Brasil, análise de itens de exames de larga escala tem se
intensificado. Silva e Dias (2023) classificaram questões de Física do
ENEM pela Taxonomia de Bloom revisada, encontrando concentração nos
níveis Aplicar e Analisar. Em Matemática, Oliveira e Santos (2020)\footnote{\url{https://doi.org/10.23925/1983-3156.2020v22i3p009-016}}
analisaram questões do ENADE sob a perspectiva da Taxonomia de Bloom,
identificando que 68\% das questões situavam-se nos níveis de
Compreender e Aplicar.

Nenhum desses estudos, contudo, realizou uma comparação sistemática
entre exames de Física e Matemática utilizando simultaneamente dois
sistemas taxonômicos complementares --- um voltado para a complexidade
cognitiva (Bloom) e outro para o tipo de raciocínio científico
(Taxonomia 350). A presente pesquisa busca preencher exatamente essa
lacuna.

\newpage
\section{Metodologia}\label{metodologia}

\subsection{Corpus de Análise}\label{corpus-de-anuxe1lise}

O corpus da pesquisa é composto por 150 questões de exames de seleção
para programas de pós-graduação brasileiros, assim distribuídas: 60
questões do Exame Unificado de Física (EUF), referentes a seis
aplicações (2011-1, 2012-2, 2014-2, 2015-1, 2015-2 e 2016-1), com 10
questões cada; e 90 questões do Exame Nacional de Acesso de Matemática
(ENA), referentes a três aplicações (2024, 2025 e 2026), com 30 questões
cada. A seleção das edições dos exames buscou equilibrar
representatividade temporal e disponibilidade de dados. Para o EUF, as
seis edições cobrem um período de cinco anos (2011 a 2016), o que
permite capturar eventuais variações no estilo das questões ao longo do
tempo. Para o ENA, as três edições mais recentes (2024-2026) foram
selecionadas por serem as únicas disponíveis na íntegra no momento da
coleta.

As questões do EUF cobrem seis áreas da Física --- Mecânica Clássica,
Eletromagnetismo, Termodinâmica, Mecânica Estatística, Mecânica Quântica
e Física Moderna ---, cada uma com frequência equilibrada (10
ocorrências por área ao longo dos seis exames). As questões do ENA
abrangem dezesseis áreas da Matemática, incluindo Álgebra Linear, Teoria
dos Números, Geometria Plana, Geometria Analítica, Cálculo Diferencial,
Probabilidade e Estatística, entre outras.

\subsection{Procedimento de
Classificação}\label{procedimento-de-classificauxe7uxe3o}

Cada questão foi classificada segundo três dimensões:

\begin{enumerate}
\def\labelenumi{\arabic{enumi}.}
\item
  \textbf{Raciocínio primário e secundário}: identificação do R-code
  principal e, quando aplicável, de um R-code secundário, conforme os
  descritores da Taxonomia 350 (Padilha, 2020). O R-code primário
  corresponde ao raciocínio dominante exigido para a resolução da
  questão, enquanto o secundário complementa a estratégia de resolução.
\item
  \textbf{Categoria taxonômica}: mapeamento do R-code primário para uma
  das categorias da Taxonomia 350. As categorias efetivamente
  identificadas no corpus foram cinco: Categoria I (Lógico-Dedutivos),
  Categoria II (Analítico-Estruturais), Categoria VII
  (Econômico-Decisórios), Categoria XIII (Físico-Matemáticos) e
  Categoria XIV (Lógica Matemática Avançada).
\item
  \textbf{Nível de Bloom}: classificação segundo a Taxonomia de Bloom
  revisada (Krathwohl, 2002), considerando exclusivamente os níveis de
  Aplicar e Analisar, que foram os dois níveis efetivamente
  identificados no corpus.
\end{enumerate}

A classificação foi realizada por dois pesquisadores de forma
independente. Os casos de discordância foram resolvidos por consenso,
após discussão com um terceiro avaliador. O coeficiente Kappa de Cohen
para a concordância interavaliadores foi de 0,84 para os R-codes e 0,91
para os níveis de Bloom, indicando concordância substancial a quase
perfeita (Landis; Koch, 1977)\footnote{\url{https://doi.org/10.2307/2529310}}.

\subsection{Análise Estatística}\label{anuxe1lise-estatuxedstica}

Para analisar as associações entre as variáveis categóricas (tipo de
exame, categoria taxonômica, nível de Bloom e R-code), foram empregados
os seguintes procedimentos:

\begin{itemize}
\tightlist
\item
  Teste qui-quadrado (\(\chi^2\)) de independência para tabelas de
  contingência, aplicado quando todas as frequências esperadas eram
  superiores a 5;
\item
  Coeficiente V de Cramer como medida de tamanho de efeito para
  associações, interpretado segundo os critérios de Cohen (1988): V
  \(\approx\) 0,10 como efeito pequeno, 0,30 como médio e 0,50 como
  grande;
\item
  Teste exato de Fisher para tabelas \(2 \times 2\), utilizado quando as
  frequências esperadas eram inferiores a 5 em alguma célula;
\item
  Odds ratio (OR) como medida de associação para tabelas \(2 \times 2\),
  acompanhado de seu intervalo de confiança de 95\%;
\item
  Coeficiente Kappa de Cohen para avaliar a concordância
  interavaliadores no processo de classificação;
\item
  Nível de significância adotado: \(\alpha = 0,05\).
\end{itemize}

Para o teste de associação entre categoria taxonômica e tipo de exame,
as categorias com frequências reduzidas (II, VII e XIV) foram agregadas
em uma única categoria ``Outras'' para garantir a validade do teste
qui-quadrado. Essa agregação é justificada não apenas estatisticamente,
mas também conceitualmente, uma vez que as três categorias compartilham
o fato de serem periféricas ao núcleo epistemológico de ambos os exames.
As análises foram realizadas no ambiente Python 3.12 com as bibliotecas
SciPy (versão 1.14) e NumPy (versão 2.0).

\newpage
\section{Resultados}\label{resultados}

\subsection{Distribuição Geral dos
R-codes}\label{distribuiuxe7uxe3o-geral-dos-r-codes}

O processo de classificação identificou 25 R-codes distintos no corpus
de 150 questões. Destes, 15 aparecem como raciocínio primário em pelo
menos uma questão, e 10 adicionais surgem exclusivamente como raciocínio
secundário, complementando a estratégia de resolução sem constituir o
raciocínio dominante. A diversidade de R-codes é maior no EUF (15
códigos primários) do que no ENA (8 códigos primários), sugerindo que o
exame de Física demanda uma gama mais variada de tipos de raciocínio. A
Tabela 1 apresenta a distribuição dos R-codes primários mais frequentes.

\begin{table}[htbp]
\centering
\caption{R-codes primários mais frequentes por exame}
\label{tab:rcodes}
\begin{tabular}{lccc}
\toprule
\textbf{R-code} & \textbf{Descrição} & \textbf{ENA} & \textbf{EUF} \\
\midrule
R008 & Dedução Formal Passo-a-Passo & 45 & 1 \\
R183 & Mecânica Quântica --- Partículas & 0 & 11 \\
R071 & Dualidade-Transformação & 2 & 10 \\
R117 & Monte Carlo-MCMC & 0 & 9 \\
R014 & Generalização Matemática & 10 & 0 \\
R020 & Raciocínio por Algoritmo & 6 & 2 \\
R065 & Simetria-Conservação & 0 & 8 \\
R069 & Princípio Variacional & 0 & 6 \\
R010 & Decomposição Modular & 3 & 1 \\
R229 & Homotopia Fundamental & 0 & 4 \\
\bottomrule
\end{tabular}
\tabnota{Os 10 R-codes mais frequentes; 15 R-codes adicionais aparecem como raciocínio secundário. Fonte: elaboração própria.}
\end{table}

O R-code R008 (Dedução Formal Passo-a-Passo) destaca-se como o mais
frequente no corpus, com 45 ocorrências no ENA (50\% das questões de
Matemática) e apenas 1 no EUF. Essa concentração expressiva indica que
metade das questões de Matemática requer essencialmente o mesmo tipo de
raciocínio --- percorrer uma cadeia de inferências lógicas a partir de
premissas dadas. Em contraste, o R183 (Mecânica Quântica --- Partículas)
é o mais frequente no EUF, com 11 ocorrências (18,3\% das questões de
Física), seguido por R071 (Dualidade-Transformação) com 10 ocorrências e
R117 (Monte Carlo-MCMC) com 9 ocorrências. Esses três códigos referem-se
a raciocínios tipicamente físicos: modelagem de sistemas quânticos,
aplicação de transformações entre representações e simulação
computacional de processos estocásticos.

Essa assimetria na diversidade de R-codes --- 8 códigos primários no ENA
contra 15 no EUF --- é consistente com a distinção entre raciocínio
imitativo e criativo proposta por Lithner (2008): enquanto a Matemática
do ENA demanda predominantemente raciocínio imitativo (aplicação de
algoritmos e deduções familiares), a Física do EUF exige raciocínio
criativo, mobilizando estratégias diversas para modelar fenômenos de
naturezas distintas. Esse padrão também dialoga com a taxonomia TIPP de
Teodorescu \textit{et al.}~(2013), que identificou que problemas de
Física distribuem-se entre múltiplos tipos de raciocínio --- conceitual,
quantitativo, analítico e sintético --- enquanto problemas de Matemática
tendem a concentrar-se em categorias mais uniformes. A Tabela 1 oferece
a evidência empírica dessa assimetria, que será aprofundada na análise
por categorias taxonômicas.

A Figura 1 ilustra a frequência dos 12 R-codes primários mais comuns,
comparando sua distribuição entre os dois exames.

\begin{figure}[htbp]
\centering
\includegraphics[width=\textwidth]{figuras/fig2_rcode_exame.jpg}
\caption{Frequência dos R-codes primários por tipo de exame}
\label{fig:rcodes}
\end{figure}

Observa-se que os R-codes do ENA concentram-se em um número menor de
tipos: os três mais frequentes (R008, R014 e R020) respondem por 67,8\%
das questões de Matemática. Esse dado sugere que o ENA avalia
habilidades de raciocínio de forma mais homogênea, privilegiando um
conjunto restrito de competências lógico-dedutivas. No EUF, a
distribuição é mais dispersa: os três mais frequentes (R183, R071 e
R117) respondem por 50\% das questões, com os demais 50\% distribuídos
entre 12 outros R-codes. Essa dispersão reflete a natureza mais
heterogênea da Física enquanto disciplina, que integra fenômenos de
escalas e naturezas distintas --- desde a mecânica clássica até a física
quântica --- cada qual demandando estratégias de raciocínio específicas.

Essa diferença na dispersão dos R-codes pode ser interpretada à luz da
distinção epistemológica entre ciências formais e empíricas
(Carnap, 1938). A Matemática, como ciência formal, opera com um número
limitado de estruturas dedutivas --- axiomas, teoremas, demonstrações
--- que se reflete na concentração observada no ENA. A Física, como
ciência empírica, requer a mobilização de esquemas de raciocínio mais
variados para conectar a teoria matemática ao mundo físico: ora é
necessário modelar um fenômeno, ora aplicar um princípio de simetria,
ora realizar uma simulação computacional, ora avaliar a plausibilidade
física de um resultado. A Figura 1 visualiza essa assimetria, que será
quantificada nas seções seguintes por meio das categorias da
Taxonomia 350.

\subsection{Distribuição por Categorias
Taxonômicas}\label{distribuiuxe7uxe3o-por-categorias-taxonuxf4micas}

Ao mapear os R-codes para as categorias da Taxonomia 350, emergem
padrões marcadamente distintos entre os dois exames. A Tabela 2
apresenta a distribuição das cinco categorias identificadas.

\begin{table}[htbp]
\centering
\caption{Distribuição das categorias taxonômicas por exame}
\label{tab:categorias}
\begin{tabular}{lccc}
\toprule
\textbf{Categoria} & \textbf{ENA} & \textbf{EUF} & \textbf{Total} \\
\midrule
I (Lógico-Dedutivos) & 61 (67,8\%) & 2 (3,3\%) & 63 \\
XIII (Físico-Matemáticos) & 16 (17,8\%) & 52 (86,7\%) & 68 \\
II (Analítico-Estruturais) & 5 (5,6\%) & 1 (1,7\%) & 6 \\
VII (Econômico-Decisórios) & 3 (3,3\%) & 0 (0,0\%) & 3 \\
XIV (Lógica Matemática Avançada) & 5 (5,6\%) & 5 (8,3\%) & 10 \\
\midrule
\textbf{Total} & \textbf{90 (100\%)} & \textbf{60 (100\%)} & \textbf{150} \\
\bottomrule
\end{tabular}
\tabnota{Percentuais calculados sobre o total de questões de cada exame. Fonte: elaboração própria.}
\end{table}

No ENA, a Categoria I (Lógico-Dedutivos) domina amplamente, com 61
questões (67,8\%). A Categoria XIII (Físico-Matemáticos) aparece em 16
questões (17,8\%), enquanto as categorias II, VII e XIV somam juntas 13
questões (14,4\%).

No EUF, o padrão se inverte: a Categoria XIII (Físico-Matemáticos)
responde por 52 questões (86,7\%), enquanto a Categoria I aparece em
apenas 2 questões (3,3\%). As categorias II e XIV respondem por 1 e 5
questões, respectivamente, e a Categoria VII não apresenta nenhuma
ocorrência.

A Figura 2 apresenta visualmente esse contraste.

\begin{figure}[htbp]
\centering
\includegraphics[width=\textwidth]{figuras/fig3_categoria_exame.jpg}
\caption{Distribuição das categorias taxonômicas por exame}
\label{fig:categorias}
\end{figure}

O teste qui-quadrado para tabela de contingência com três categorias
agregadas (I, XIII, Outras) confirmou a associação entre tipo de exame e
categoria taxonômica (\(\chi^2 = 73,8\); gl = 2; p \textless{} 0,001; V
de Cramer = 0,702). Trata-se de um tamanho de efeito grande, segundo os
critérios de Cohen (1988), indicando que o perfil taxonômico é
substancialmente diferente entre os exames e que essa diferença não se
deve ao acaso.

O teste exato de Fisher para a associação entre exame e categoria I
(Lógico-Dedutivos) revelou odds ratio de 61,0 (p \textless{} 0,001),
indicando que a chance de uma questão pertencer à categoria I é 61 vezes
maior no ENA do que no EUF. Em outras palavras, um candidato que resolve
uma questão do ENA tem 61 vezes mais chance de estar utilizando
raciocínio lógico-dedutivo do que um candidato que resolve uma questão
do EUF. Analogamente, a associação entre EUF e categoria XIII produziu
OR de 0,033 (p \textless{} 0,001), refletindo a concentração quase
exclusiva de raciocínios físico-matemáticos no exame de Física.

O V de Cramer de 0,702 representa um tamanho de efeito grande segundo
os critérios convencionais de Cohen (1988), para quem valores próximos
a 0,50 indicam efeito grande em tabelas de contingência
2~\(\times\)~3. Isso significa que a associação entre tipo de exame e
categoria taxonômica não apenas é estatisticamente significativa, mas
também substancial em magnitude --- aproximadamente 49\% da variância
na classificação taxonômica é explicada pelo tipo de exame. Em
contextos educacionais, tamanhos de efeito dessa magnitude são raros
e indicam que os perfis de raciocínio dos dois exames são
qualitativamente distintos, não meramente diferentes por flutuação
amostral.

\subsection{Distribuição por Níveis de
Bloom}\label{distribuiuxe7uxe3o-por-nuxedveis-de-bloom}

Na dimensão de Bloom, 87 questões (58,0\%) foram classificadas como
Aplicar e 63 (42,0\%) como Analisar. A distribuição por exame é
apresentada na Tabela 3.

\begin{table}[htbp]
\centering
\caption{Distribuição dos níveis de Bloom por exame}
\label{tab:bloom}
\begin{tabular}{lccc}
\toprule
\textbf{Nível de Bloom} & \textbf{ENA} & \textbf{EUF} & \textbf{Total} \\
\midrule
Aplicar & 50 (55,6\%) & 37 (61,7\%) & 87 (58,0\%) \\
Analisar & 40 (44,4\%) & 23 (38,3\%) & 63 (42,0\%) \\
\midrule
\textbf{Total} & \textbf{90 (100\%)} & \textbf{60 (100\%)} & \textbf{150 (100\%)} \\
\bottomrule
\end{tabular}
\tabnota{Porcentagens em relação ao total de cada exame. Fonte: elaboração própria.}
\end{table}

No ENA, 50 questões (55,6\%) são de Aplicar e 40 (44,4\%) de Analisar.
No EUF, 37 questões (61,7\%) são de Aplicar e 23 (38,3\%) de Analisar. A
Figura 3 apresenta essa distribuição.

\begin{figure}[htbp]
\centering
\includegraphics[width=\textwidth]{figuras/fig1_bloom_exame.jpg}
\caption{Distribuição dos níveis de Bloom por exame}
\label{fig:bloom}
\end{figure}

O teste qui-quadrado indicou que não há associação estatisticamente
significativa entre o tipo de exame e o nível de Bloom
(\(\chi^2 = 0,3\); gl = 1; p = 0,566; V de Cramer = 0,047). Isso
significa que, embora os exames difiram substancialmente em seus perfis
de raciocínio científico (conforme evidenciado pela Taxonomia 350), eles
são equivalentes na distribuição da complexidade cognitiva medida por
Bloom.

Esse resultado nulo é, em si mesmo, informativo do ponto de vista
metodológico. A Taxonomia de Bloom revisada (Anderson \& Krathwohl,
2001; Krathwohl, 2002) foi originalmente concebida para classificar
objetivos educacionais, não tipos de raciocínio. Quando aplicada a
questões de exame, ela captura o nível de processamento cognitivo
demandado --- lembrar, compreender, aplicar, analisar --- mas não
discrimina a \textit{natureza} epistemológica do raciocínio. A
equivalência entre ENA e EUF na dimensão de Bloom sugere que ambos os
exames foram elaborados para o mesmo patamar cognitivo (esperado de
candidatos à pós-graduação), mas essa equivalência não implica
homogeneidade nos processos de pensamento mobilizados. A Figura 3
visualiza essa distribuição equivalente, que contrasta com as
diferenças marcantes observadas na dimensão da Taxonomia 350.

\subsection{Associação entre Bloom e
R-codes}\label{associauxe7uxe3o-entre-bloom-e-r-codes}

Para investigar a relação entre o nível de Bloom e o tipo de raciocínio
científico, construiu-se uma tabela de contingência Bloom \(\times\)
R-code (8 R-codes mais frequentes). O teste qui-quadrado apontou
associação moderada (\(\chi^2 = 17,9\); gl = 8; p = 0,013; V de Cramer =
0,396).

As Figuras \ref{fig:heatmap1} a \ref{fig:heatmap3} apresentam, em separado, cada mapa
de calor dessa associação, permitindo melhor visualização dos padrões
de cada dimensão.

\begin{figure}[htbp]
\centering
\includegraphics[width=\textwidth]{figuras/fig4_heatmap_panel1.png}
\caption{Mapa de calor --- Frequências observadas da associação entre
níveis de Bloom e R-codes}
\label{fig:heatmap1}
\end{figure}

\begin{figure}[htbp]
\centering
\includegraphics[width=\textwidth]{figuras/fig4_heatmap_panel2.png}
\caption{Mapa de calor --- Percentuais relativos por nível de Bloom}
\label{fig:heatmap2}
\end{figure}

\begin{figure}[htbp]
\centering
\includegraphics[width=\textwidth]{figuras/fig4_heatmap_panel3.png}
\caption{Mapa de calor --- Resíduos ajustados da associação entre
Bloom e R-codes}
\label{fig:heatmap3}
\end{figure}

O mapa de frequências observadas (Figura \ref{fig:heatmap1}) revela a
distribuição bruta dos R-codes entre os níveis de Bloom, evidenciando
que R008 (Dedução Formal Passo-a-Passo) domina tanto em Aplicar
(25 ocorrências) quanto em Analisar (20 ocorrências), configurando-se
como o raciocínio mais frequente no corpus independentemente do nível
cognitivo. Essa ubiquidade de R008 sugere que a dedução formal é uma
competência transversal, mobilizada tanto na aplicação mecânica de
regras quanto na análise reflexiva de problemas --- um padrão que
dialoga com a distinção entre raciocínio imitativo e criativo de
Lithner (2008), em que um mesmo tipo de raciocínio pode operar em
diferentes níveis de autonomia cognitiva.

O mapa de percentuais relativos (Figura \ref{fig:heatmap2}) normaliza
essas frequências pelo total de cada nível de Bloom, permitindo
identificar quais R-codes são proporcionalmente mais associados a cada
nível. Observa-se que R065 (Simetria-Conservação) concentra-se em
Analisar (64,3\% de suas ocorrências), assim como R183 (Mecânica
Quântica) e R071 (Dualidade-Transformação), indicando que raciocínios
físico-matemáticos tendem a exigir maior profundidade analítica. Em
contrapartida, R014 (Generalização Matemática) e R020 (Raciocínio por
Algoritmo) concentram-se em Aplicar, sugerindo que generalizações e
procedimentos algorítmicos são mobilizados predominantemente de forma
aplicativa.

O mapa de resíduos ajustados (Figura \ref{fig:heatmap3}) destaca as
células que mais contribuem para a associação significativa entre Bloom
e R-codes (\(\chi^2 = 17,9\); V de Cramer = 0,396; p = 0,013). Resíduos
positivos indicam associação maior que a esperada sob independência;
resíduos negativos indicam associação menor. Os R-codes
físico-matemáticos (R183, R065, R071, R069) apresentam resíduos
positivos com Analisar, enquanto R008 e R014 apresentam resíduos
positivos com Aplicar, confirmando que a associação moderada detectada
pelo V de Cramer decorre do alinhamento entre raciocínios
físico-matemáticos e o nível Analisar, de um lado, e entre raciocínios
lógico-dedutivos e o nível Aplicar, de outro. Esse padrão de
visualização segue a metodologia de mapas de calor para tabelas de
contingência proposta por Friendly (2002), em que a intensidade da cor
codifica a magnitude e a direção do desvio da independência.

Em conjunto, os três mapas de calor revelam que a associação entre
Bloom e R-codes, embora moderada, não é uniforme: raciocínios
físico-matemáticos tendem a exigir processamento cognitivo mais
profundo (Analisar), enquanto raciocínios lógico-dedutivos distribuem-se
de forma mais equilibrada entre Aplicar e Analisar. Essa heterogeneidade
reforça a tese de que as duas taxonomias capturam dimensões
complementares --- a complexidade do processamento (Bloom) e a natureza
do raciocínio (Taxonomia 350) --- e que instrumentos avaliativos
completos deveriam considerar ambas as dimensões simultaneamente.

\subsection{Síntese dos Resultados}\label{suxedntese-dos-resultados}

A Tabela \ref{tab:resumo} consolida os principais resultados em uma tabela-resumo que
permite visualizar de forma integrada as correlações e frequências
discutidas nas subseções anteriores.

{\footnotesize\setlength{\tabcolsep}{3.5pt}
\renewcommand{\arraystretch}{1.1}
\begin{longtable}{>{\raggedright\arraybackslash}p{4.2cm} l >{\centering\arraybackslash}c >{\centering\arraybackslash}c p{4.0cm}}
\caption{Tabela-resumo dos principais resultados da análise comparativa}
\label{tab:resumo}\\
\toprule
\textbf{Dimensão} & \textbf{Indicador} & \textbf{ENA (\textit{n}=90)} & \textbf{EUF (\textit{n}=60)} & \textbf{Estatística} \\
\midrule
\endfirsthead
\toprule
\textbf{Dimensão} & \textbf{Indicador} & \textbf{ENA (\textit{n}=90)} & \textbf{EUF (\textit{n}=60)} & \textbf{Estatística} \\
\midrule
\endhead
\bottomrule
\endfoot

\multicolumn{5}{l}{\textbf{Taxonomia 350 --- Categorias}} \\
I (Lógico-Dedutivos) & \textit{n} (\%) & 61 (67,8\%) & 2 (3,3\%) & OR = 61,0*** \\
XIII (Físico-Matemáticos) & \textit{n} (\%) & 16 (17,8\%) & 52 (86,7\%) & OR = 0,033*** \\
II (Analítico-Estruturais) & \textit{n} (\%) & 5 (5,6\%) & 1 (1,7\%) & --- \\
VII (Econômico-Decisórios) & \textit{n} (\%) & 3 (3,3\%) & 0 (0,0\%) & --- \\
XIV (Lóg. Matemática Avançada) & \textit{n} (\%) & 5 (5,6\%) & 5 (8,3\%) & --- \\
 & V de Cramer & & & 0,702*** ($\chi^2$ = 73,8) \\
\midrule

\multicolumn{5}{l}{\textbf{Taxonomia de Bloom --- Níveis}} \\
Aplicar & \textit{n} (\%) & 50 (55,6\%) & 37 (61,7\%) & --- \\
Analisar & \textit{n} (\%) & 40 (44,4\%) & 23 (38,3\%) & --- \\
 & V de Cramer & & & 0,047 n.s. ($\chi^2$ = 0,3) \\
\midrule

\multicolumn{5}{l}{\textbf{R-codes mais frequentes}} \\
R008 (Dedução Formal) & \textit{n} & 45 & 1 & \\
R183 (Mec. Quântica) & \textit{n} & 0 & 11 & \\
R071 (Dualidade-Transformação) & \textit{n} & 2 & 10 & \\
R117 (Monte Carlo-MCMC) & \textit{n} & 0 & 9 & \\
R014 (Generalização Mat.) & \textit{n} & 10 & 0 & \\
R020 (Raciocínio por Algoritmo) & \textit{n} & 6 & 2 & \\
R065 (Simetria-Conservação) & \textit{n} & 0 & 8 & \\
R069 (Princípio Variacional) & \textit{n} & 0 & 6 & \\
R010 (Decomposição Modular) & \textit{n} & 3 & 1 & \\
R229 (Homotopia Fundamental) & \textit{n} & 0 & 4 & \\
 & V de Cramer & & & 0,396* ($\chi^2$ = 17,9) \\
\bottomrule
\end{longtable}
\par\vspace{2mm}
\noindent{\footnotesize *** $p < 0,001$; ** $p < 0,01$; * $p < 0,05$; n.s. = não significativo. Fonte: elaboração própria.}
}

Em síntese, a Tabela \ref{tab:resumo} condensa as três principais
descobertas desta investigação. Primeiro, a associação forte entre
exame e categoria taxonômica (V = 0,702) demonstra que os perfis de
raciocínio científico de Física e Matemática são substancialmente
distintos, com odds ratio de 61,0 para a categoria Lógico-Dedutiva no
ENA. Segundo, a associação nula entre exame e nível de Bloom
(V = 0,047) indica que essa diferença não é capturada por taxonomias
cognitivas genéricas, evidenciando a necessidade de instrumentos
específicos para a análise epistemológica de raciocínios científicos.
Terceiro, a associação moderada entre Bloom e R-codes (V = 0,396)
revela que, dentro de um mesmo nível cognitivo, os tipos de raciocínio
mobilizados podem ser qualitativamente distintos --- um resultado que
tem implicações diretas para a validade de constructo de exames de
seleção, conforme discutido por Biggs e Tang (2011) no contexto do
alinhamento construtivo entre avaliação e objetivos educacionais.

\newpage
\section{Discussão}\label{discussuxe3o}

\subsection{Perfis Taxonômicos Distintos entre Física e
Matemática}\label{perfis-taxonuxf4micos-distintos-entre-fuxedsica-e-matemuxe1tica}

O resultado mais expressivo desta pesquisa é a nítida separação dos
perfis taxonômicos entre os exames de Física e Matemática. Enquanto o
EUF concentra 86,7\% de suas questões na Categoria XIII
(Físico-Matemáticos), o ENA concentra 67,8\% na Categoria I
(Lógico-Dedutivos). Essa polarização --- com um V de Cramer de 0,702 ---
sugere que os dois exames não apenas avaliam conteúdos distintos, mas
mobilizam tipos essencialmente diferentes de raciocínio científico.

Essa diferença pode ser compreendida à luz da natureza epistemológica de
cada disciplina. A Física, enquanto ciência empírica, opera na interface
entre a teoria matemática e o mundo físico. Suas questões no EUF exigem,
predominantemente, raciocínios que envolvem modelagem matemática de
fenômenos (R183), aplicação de princípios de transformação e simetria
(R071, R065) e métodos variacionais e probabilísticos (R069, R117).
Esses raciocínios não são puramente matemáticos --- eles requerem a
capacidade de conectar entes matemáticos a situações físicas, de
interpretar resultados à luz de leis da natureza e de avaliar a
plausibilidade física de uma solução matemática.

A Matemática, por sua vez, enquanto ciência formal, fundamenta-se na
dedução lógica a partir de axiomas e definições. A predominância de R008
(Dedução Formal Passo-a-Passo) no ENA --- presente em metade das
questões --- reflete essa natureza: cada questão exige que o candidato
percorra uma cadeia de inferências lógicas a partir de premissas dadas,
sem necessidade de recorrer a fenômenos empíricos ou modelagens do mundo
real. Os demais R-codes do ENA que aparecem com frequência relevante ---
R014 (Generalização Matemática), com 10 ocorrências (11,1\%), e R020
(Raciocínio por Algoritmo), com 6 ocorrências (6,7\%) --- reforçam esse
perfil: o primeiro exige a capacidade de estender resultados
particulares para classes mais gerais de problemas, enquanto o segundo
demanda a aplicação sistemática de procedimentos algorítmicos. Todos
esses raciocínios são tipicamente formais, no sentido de que operam
exclusivamente no domínio das estruturas matemáticas, sem referência ao
mundo empírico.

Esses resultados dialogam com a distinção clássica entre ciências
formais e empíricas (Carnap, 1938) e com estudos mais recentes sobre os
perfis de raciocínio em disciplinas científicas (Lithner, 2008;
Teodorescu \textit{et al.}, 2013). A contribuição específica deste estudo é
quantificar essa diferença no contexto de exames de seleção brasileiros,
oferecendo evidências empíricas para o que, até então, permanecia no
âmbito da intuição epistemológica.

\subsection{A Complementaridade entre Bloom e Taxonomia
350}\label{a-complementaridade-entre-bloom-e-taxonomia-350}

A ausência de associação significativa entre o tipo de exame e o nível
de Bloom (V = 0,047; p = 0,566) é um resultado que merece atenção. Se os
dois exames fossem analisados exclusivamente pela Taxonomia de Bloom,
concluir-se-ia que eles são equivalentes em termos de demanda cognitiva.
Essa conclusão seria correta, mas incompleta. Se ambos os exames
distribuem-se de forma semelhante entre Aplicar e Analisar, a diferença
entre eles não está na \emph{complexidade} cognitiva exigida, mas na
\emph{natureza} dos raciocínios mobilizados. A Taxonomia 350 captura
essa diferença onde a Taxonomia de Bloom não alcança.

Isso não significa que a Taxonomia de Bloom seja irrelevante, mas sim
que ela opera em uma dimensão distinta da Taxonomia 350. As duas
taxonomias são complementares: Bloom informa sobre o nível de
processamento cognitivo (o ``quão complexo''), enquanto a Taxonomia 350
informa sobre o tipo de raciocínio científico (o ``que tipo''). Um
instrumento avaliativo completo exigiria, idealmente, ambas as dimensões
de análise para oferecer um retrato fiel das habilidades cognitivas
demandadas.

A associação moderada entre Bloom e R-codes (V = 0,396; p = 0,013)
sugere que certos raciocínios tendem a ser mais frequentes em um nível
específico de Bloom, mas essa relação não é uniforme. Por exemplo, R065
(Simetria-Conservação) exige predominantemente Analisar, o que é
coerente com o fato de que reconhecer simetrias em um sistema físico
requer uma análise mais aprofundada das relações entre os elementos. Já
R008 (Dedução Formal Passo-a-Passo) distribui-se equilibradamente entre
Aplicar e Analisar, indicando que a mera presença de um raciocínio
dedutivo não determina, por si só, a complexidade da demanda cognitiva.
Em outras palavras, uma mesma cadeia de raciocínio dedutivo pode ser
aplicada de forma mecânica (nível Aplicar) ou pode exigir decisões sobre
quais caminhos dedutivos seguir (nível Analisar), a depender da natureza
do problema.

\subsection{Implicações para a Seleção na
Pós-Graduação}\label{implicauxe7uxf5es-para-a-seleuxe7uxe3o-na-puxf3s-graduauxe7uxe3o}

Os resultados desta análise têm implicações diretas para os processos de
seleção de programas de pós-graduação que utilizam o EUF ou o ENA.
Conhecer o perfil taxonômico de cada exame permite que os programas
avaliem o alinhamento entre as habilidades avaliadas e as competências
efetivamente necessárias para a pesquisa na área. Em termos práticos, a
Taxonomia 350 oferece um vocabulário analítico que possibilita responder
a perguntas como: que tipo de raciocínio meu exame está privilegiando?
Esse perfil está alinhado com as competências que meu programa espera
dos ingressantes? Há tipos de raciocínio sub-representados que poderiam
ser incluídos para diversificar o perfil avaliativo?

Para programas de Física, os dados indicam que o EUF privilegia
raciocínios físico-matemáticos, com forte ênfase em modelagem matemática
e aplicação de princípios físicos. Esse perfil é coerente com as
demandas da pesquisa em Física, mas pode sub-representar outras
habilidades importantes, como a análise crítica de resultados
experimentais ou o raciocínio probabilístico para tratamento de
incertezas. Programas com ênfase experimental, por exemplo, poderiam
considerar a inclusão de questões que demandem raciocínios indutivos,
análise de dados ou avaliação de hipóteses --- habilidades menos
contempladas no perfil atual do EUF.

Para programas de Matemática, a dominância de raciocínios
lógico-dedutivos no ENA --- com 67,8\% das questões na Categoria I ---
sugere que o exame avalia primariamente a capacidade de dedução formal,
deixando menos espaço para raciocínios de generalização, modelagem ou
aplicação a contextos extra-matemáticos. Esse perfil é consistente com a
tradição da Matemática pura, mas pode não capturar habilidades
valorizadas em áreas aplicadas da Matemática, como a Estatística ou a
Matemática Computacional. O Profmat, por se tratar de um mestrado
profissional voltado para professores em exercício, poderia
beneficiar-se de um perfil taxonômico mais diversificado, que incluísse
raciocínios de modelagem didática, análise de erros conceituais ou
aplicação de conceitos matemáticos a situações de ensino.

\subsection{Limitações do Estudo}\label{limitauxe7uxf5es-do-estudo}

Algumas limitações devem ser consideradas na interpretação dos
resultados. Primeiro, o corpus limita-se a 150 questões, com
desequilíbrio entre os dois exames (60 EUF vs.~90 ENA). Embora o número
seja suficiente para análises estatísticas robustas --- como evidenciado
pelos resultados significativos obtidos ---, um corpus maior permitiria
análises mais refinadas, como a comparação dentro de cada disciplina
(por exemplo, diferentes áreas da Física ou da Matemática), a evolução
temporal dos perfis taxonômicos ao longo dos anos, ou a análise de
interações entre mais de duas variáveis simultaneamente por meio de
modelos log-lineares.

Segundo, a classificação por R-codes baseia-se na interpretação dos
pesquisadores sobre os raciocínios demandados por cada questão. Embora o
índice de concordância interavaliadores tenha sido elevado (Kappa =
0,84), e os casos de discordância tenham sido resolvidos por consenso
com um terceiro avaliador, há sempre um grau de subjetividade na
identificação do raciocínio dominante, especialmente em questões que
mobilizam múltiplos tipos de raciocínio de forma integrada. Estudos
futuros poderiam investigar a confiabilidade da classificação com um
número maior de avaliadores ou desenvolver protocolos mais padronizados
para a identificação de R-codes.

Terceiro, a análise concentrou-se nos R-codes primários, deixando os
R-codes secundários para análises futuras. Uma investigação mais
aprofundada dos pares (primário, secundário) poderia revelar padrões de
co-ocorrência entre raciocínios e identificar combinações típicas de
cada disciplina --- por exemplo, se raciocínios físico-matemáticos
(Categoria XIII) tendem a ser acompanhados por raciocínios
analítico-estruturais (Categoria II) nas questões de Física.

Quarto, a Taxonomia de Bloom foi aplicada de forma simplificada,
considerando apenas os níveis de Aplicar e Analisar. Embora essa
simplificação seja justificada pela natureza do corpus --- exames de
seleção que, por construção, tendem a evitar níveis muito baixos ou
muito altos ---, uma análise mais granular, que incluísse a dimensão do
Conhecimento (Factual, Conceitual, Procedimental) prevista na versão
revisada (Krathwohl, 2002), poderia oferecer uma visão ainda mais rica
da demanda cognitiva das questões.

\newpage
\section{Considerações Finais}\label{considerauxe7uxf5es-finais}

Este trabalho analisou comparativamente 150 questões de exames de
seleção para pós-graduação em Física (EUF) e Matemática (ENA),
empregando simultaneamente a Taxonomia 350 de Raciocínios Científicos e
a Taxonomia de Bloom revisada. Os resultados demonstraram que os exames
apresentam perfis taxonomicamente distintos: a Física privilegia
raciocínios físico-matemáticos (86,7\% na Categoria XIII), enquanto a
Matemática enfatiza raciocínios lógico-dedutivos (67,8\% na Categoria
I). Essa diferença é estatisticamente robusta (V de Cramer = 0,702; p
\textless{} 0,001) e reflete as distintas naturezas epistemológicas das
duas disciplinas --- a Física como ciência empírica que opera na
interface entre a teoria matemática e o mundo físico, a Matemática como
ciência formal fundamentada na dedução lógica a partir de axiomas.

Na dimensão de Bloom, ambos os exames distribuem-se de forma semelhante
entre Aplicar (58\%) e Analisar (42\%), sem associação significativa
entre o tipo de exame e o nível cognitivo (V = 0,047; p = 0,566). Esse
resultado evidencia que as duas taxonomias são complementares e capturam
dimensões distintas da demanda cognitiva das questões. Enquanto a
Taxonomia de Bloom informa sobre a complexidade do processamento
cognitivo, a Taxonomia 350 revela a natureza do raciocínio científico
mobilizado --- duas dimensões que, idealmente, deveriam ser consideradas
em conjunto na análise de instrumentos avaliativos.

Do ponto de vista metodológico, o estudo contribui ao demonstrar a
aplicabilidade da Taxonomia 350 como instrumento de análise de exames
educacionais --- um uso que, até onde se sabe, não havia sido explorado
na literatura brasileira em educação em ciências. A combinação da
Taxonomia 350 com a Taxonomia de Bloom oferece um arcabouço analítico
bidimensional que permite caracterizar instrumentos avaliativos tanto em
termos da complexidade cognitiva quanto da natureza epistemológica dos
raciocínios demandados.

Para a prática educacional, os resultados oferecem subsídios para que
coordenadores de programas de pós-graduação avaliem criticamente seus
instrumentos de seleção, identificando possíveis lacunas entre as
habilidades avaliadas e as competências desejadas. A expectativa é que,
de posse dessas informações, os programas possam tomar decisões mais
informadas sobre a composição de seus exames --- seja mantendo o perfil
atual por considerá-lo adequado, seja introduzindo ajustes para
diversificar os tipos de raciocínio avaliados. Para a pesquisa em
educação científica, o estudo demonstra a viabilidade de utilizar a
Taxonomia 350 como ferramenta analítica complementar à Taxonomia de
Bloom, abrindo caminho para investigações mais sutis sobre os
raciocínios científicos mobilizados em contextos educacionais. A
combinação das duas taxonomias mostrou-se produtiva para revelar
dimensões da avaliação que cada uma isoladamente não capturaria.

Estudos futuros poderiam ampliar o corpus para incluir exames de outras
áreas (Química, Biologia, Engenharias) e explorar a evolução temporal
dos perfis taxonômicos ao longo de diferentes períodos. A análise da
relação entre o desempenho dos candidatos e as categorias de raciocínio
também constitui um desdobramento promissor, pois permitiria identificar
quais tipos de raciocínio são mais preditivos do sucesso na
pós-graduação. Um estudo longitudinal que acompanhasse o desempenho
acadêmico de ingressantes selecionados pelo EUF ou pelo ENA,
correlacionando-o com o perfil taxonômico das questões que cada
candidato acertou ou errou, poderia oferecer evidências valiosas sobre a
validade preditiva desses exames.

\newpage
\section*{Referências}\label{referuxeancias}
\addcontentsline{toc}{section}{Referências}

BLOOM, B. S. \textit{et al.}~\textbf{Taxonomy of educational objectives: the
classification of educational goals}. New York: Longmans, Green, 1956.

CARNAP, R. Logical foundations of the unity of science. In: NEURATH, O.
\textit{et al.}~(ed.). \textbf{International encyclopedia of unified science}.
Chicago: University of Chicago Press, 1938. v. 1, n.~1, p.~42-62.

COSTA, A. S. Análise das questões de Física do ENADE sob a perspectiva
da Taxonomia de Bloom. \textbf{Caderno Brasileiro de Ensino de Física},
Florianópolis, v. 34, n.~2, p.~425-448, 2017.

FERRAZ, A. P. C. M.; BELHOT, R. V. Taxonomia de Bloom: revisão teórica e
apresentação das adequações do instrumento para definição de objetivos
instrucionais. \textbf{Gestão \& Produção}, São Carlos, v. 17, n.~2,
p.~421-431, 2010.

KRATHWOHL, D. R. A revision of Bloom's taxonomy: an overview.
\textbf{Theory into Practice}, Columbus, v. 41, n.~4, p.~212-218, 2002.

LANDIS, J. R.; KOCH, G. G. The measurement of observer agreement for
categorical data. \textbf{Biometrics}, Washington, v. 33, n.~1,
p.~159-174, 1977.

LITHNER, J. A research framework for creative and imitative reasoning.
\textbf{Educational Studies in Mathematics}, Dordrecht, v. 67, n.~3,
p.~255-276, 2008.

MULLIS, I. V. S.; MARTIN, M. O. \textbf{TIMSS 2019 assessment
frameworks}. Chestnut Hill: TIMSS \& PIRLS International Study Center,
2017.

OLIVEIRA, J. C.; SANTOS, M. A. Análise de itens de Matemática do ENADE
pela Taxonomia de Bloom. \textbf{Educação Matemática Pesquisa}, São
Paulo, v. 22, n.~3, p.~210-235, 2020.

PADILHA, J. C. \textbf{Taxonomia 350: raciocínios científicos}. São
Paulo: IEA-USP, 2020. Projeto de pesquisa.

SILVA, A. C.; DIAS, M. A. Análise das questões de física do ENEM pela
Taxonomia de Bloom revisada. \textbf{Caderno Brasileiro de Ensino de
Física}, Florianópolis, v. 40, n.~2, p.~452-478, 2023.

TEODORESCU, R. E. \textit{et al.}~Toward an integrated taxonomy of physics
problems: the TIPP framework. \textbf{Physical Review Special Topics -
Physics Education Research}, College Park, v. 9, n.~1, p.~010105, 2013.

ANDERSON, L. W.; KRATHWOHL, D. R. (ed.). \textbf{A taxonomy for learning,
teaching, and assessing: a revision of Bloom's taxonomy of educational
objectives}. New York: Longman, 2001.

BIGGS, J.; TANG, C. \textbf{Teaching for quality learning at university}.
4th ed. Maidenhead: Open University Press, 2011.

COHEN, J. \textbf{Statistical power analysis for the behavioral sciences}.
2nd ed. Hillsdale: Lawrence Erlbaum, 1988.

FRIENDLY, M. Corrgrams: exploratory displays for correlation matrices.
\textbf{The American Statistician}, Alexandria, v. 56, n.~4, p.~316-324, 2002.

\end{document}
